\documentclass{amsart}
\usepackage{amsmath,amssymb,amsthm,hyperref,mathtools}

\newtheorem{theorem}{Theorem}
\newtheorem{lemma}{Lemma}
\newtheorem{proposition}{Proposition}
\newtheorem{corollary}{Corollary}
\theoremstyle{definition}
\newtheorem{definition}{Definition}
\newtheorem{remark}{Remark}
\newtheorem{example}{Example}

\DeclareMathOperator{\uce}{uce}
\DeclareMathOperator{\HC}{HC}
\DeclareMathOperator{\im}{im}
\DeclareMathOperator{\Hom}{Hom}
\DeclareMathOperator{\id}{id}
\DeclareMathOperator{\Aut}{Aut}
\newcommand{\g}{\mathfrak{g}}
\newcommand{\h}{\mathfrak{h}}
\newcommand{\nn}{\mathfrak{n}}
\newcommand{\Om}{\Omega}
\newcommand{\Gm}{\Gamma}
\newcommand{\ra}{\rightarrow}

\begin{document}

\title{Universal central extensions of equivariant map algebras}
\maketitle

Throughout, $k$ is a field of characteristic $0$, $\g$ is a finite-dimensional
Lie algebra over $k$, $A$ is a unital commutative associative $k$-algebra, and
$\Gm$ is a finite group acting by Lie algebra automorphisms on $\g$ and by algebra
automorphisms on $A$. We let $\Gm$ act diagonally on the current algebra
$L:=\g\ot A$, and set
\[
   M=L^{\Gm}=(\g\ot A)^{\Gm}=\{u\in\g\ot A:\gamma\cdot u=u\ \text{for all }\gamma\in\Gm\}.
\]
We assume $M$ is perfect.  All tensor products, $\Hom$'s, exterior and symmetric
powers are over $k$ unless decorated otherwise.

Because $\operatorname{char}k=0$ and $\Gm$ is finite, the averaging idempotent
\[
   e=\frac{1}{|\Gm|}\sum_{\gamma\in\Gm}\gamma\in k[\Gm]
\]
acts on every $k[\Gm]$-module $V$ as a projector with image $V^{\Gm}$, giving a
canonical $\Gm$-stable decomposition $V=V^{\Gm}\oplus(1-e)V$.  Hence the functor
$(-)^{\Gm}$ is \emph{exact} on $k[\Gm]$-modules (it is the direct summand $e$ of
the identity functor), and $e$ commutes with every $\Gm$-equivariant linear map.
We use these two facts constantly.

\paragraph{Summary of results.}  Part (a) is settled completely and unconditionally
(Theorem~\ref{thm:partA}): $\uce(M)=\Lambda^2M/J_M$ with centre $H_2(M)$.  For part
(b) we prove (Proposition~\ref{prop:snake}) that the natural comparison map $\Phi$
is an isomorphism \emph{iff} an explicit linear map
$\phi\colon H_2(M)\to H_2(\g\ot A)^{\Gm}$ is.  We then determine $\phi$:
it is an isomorphism for $\g$ semisimple (Theorem~\ref{thm:partB}(4); injectivity
proved in full, surjectivity via Kassel's computation and a Casimir descent), and more generally whenever the
isotypic complement $(1-e)L$ is an ideal (Theorem~\ref{thm:partB}(3)); it can fail
to be surjective when $\g$ is merely perfect with $H_2(\g)\ne0$
(Theorem~\ref{thm:partB}(5), Lemma~\ref{lem:strict}).

\bigskip
\section{Universal central extensions: the machinery}

\begin{definition}\label{def:uce}
For a Lie algebra $L$ over $k$, a \emph{central extension} is a surjective Lie
homomorphism $\pi\colon \widehat L\ra L$ with $\ker\pi\subseteq Z(\widehat L)$. It
is \emph{universal} if for every central extension $\pi'\colon E\ra L$ there is a
unique Lie homomorphism $\varphi\colon\widehat L\ra E$ with
$\pi'\circ\varphi=\pi$.
\end{definition}

\begin{theorem}[Garland--Kassel]\label{thm:uce-exists}
Let $L$ be a perfect Lie algebra over $k$.  Put
\[
   \uce(L)=\Lambda^2L/J,\quad
   J=\operatorname{span}_k\{[x,y]\wedge z+[y,z]\wedge x+[z,x]\wedge y:x,y,z\in L\},
\]
write $\langle x,y\rangle$ for the image of $x\wedge y$, define the bracket
$[\langle x,y\rangle,\langle u,v\rangle]=\langle[x,y],[u,v]\rangle$ and the map
$\pi\colon\uce(L)\ra L$, $\langle x,y\rangle\mapsto[x,y]$.  Then $\pi$ is the
universal central extension of $L$, the algebra $\uce(L)$ is perfect, and
\[
   \ker\pi=Z(\uce(L))\cong H_2(L):=H_2(L;k),
\]
the second Chevalley--Eilenberg homology with trivial coefficients.  Concretely,
with the boundary maps
\[
   \Lambda^3L\xrightarrow{d_3}\Lambda^2L\xrightarrow{d_2}L,\quad
   d_2(x\wedge y)=[x,y],\
   d_3(x\wedge y\wedge z)=[x,y]\wedge z-[x,z]\wedge y+[y,z]\wedge x,
\]
one has $H_2(L)=\ker d_2/\im d_3=\ker d_2/J$, equal to $\ker\pi\subseteq\Lambda^2L/J$.
\end{theorem}

\begin{proof}
Classical (Weibel, \emph{Introduction to Homological Algebra}, \S7.9;
Kassel--Loday, \emph{Ann. Inst. Fourier} 32 (1982); Garland, \emph{Publ. IHÉS} 52
(1980)).  We record the points used.

\emph{Bracket well defined.}  The bilinear map
$(x\wedge y,u\wedge v)\mapsto[x,y]\wedge[u,v]$ on $\Lambda^2L$ kills $J$ in each
slot: for $j=[x,y]\wedge z+[y,z]\wedge x+[z,x]\wedge y$,
\[
   [j,u\wedge v]\mapsto\bigl([[x,y],z]+[[y,z],x]+[[z,x],y]\bigr)\wedge[u,v]=0
\]
by the Jacobi identity, and symmetrically; Jacobi for $\uce(L)$ reduces to Jacobi
for $L$.

\emph{$\pi$ central, $\uce(L)$ perfect.}  $\pi$ is a surjective homomorphism since
$L=[L,L]$; if $[x,y]=0$ then $[\langle x,y\rangle,\langle u,v\rangle]=
\langle0,[u,v]\rangle=0$, so $\ker\pi$ is central.  A generator $\langle
x,y\rangle$ with $y=\sum_i[c_i,d_i]$ equals $\sum_i[\langle x,c_i\rangle,\langle
e_i,f_i\rangle]$ ($[e_i,f_i]=d_i$), so $\uce(L)=[\uce(L),\uce(L)]$.

\emph{Universality.}  For a central extension $\pi'\colon E\ra L$ with $k$-linear
section $s$, set $\varphi\langle x,y\rangle=[s(x),s(y)]$; centrality of $\ker\pi'$
makes this independent of $s$ and kills $J$, giving a Lie homomorphism with
$\pi'\varphi=\pi$.  Two lifts of $\pi$ differ by a central map vanishing on
$\uce(L)=[\uce(L),\uce(L)]$, hence coincide.

\emph{Homology.}  For a perfect central extension, $Z(\widehat L)=\ker\pi$, and
$\ker\pi=\ker d_2/\im d_3=H_2(L)$ by the displayed Chevalley--Eilenberg complex.
\end{proof}

\begin{corollary}[Functoriality and induced group actions]\label{cor:funct}
$\uce$ is a functor on perfect Lie algebras: $f\colon L\ra L'$ induces the unique
lift $\uce(f)\colon\uce(L)\ra\uce(L')$,
$\langle x,y\rangle\mapsto\langle f(x),f(y)\rangle$, with
$\pi'\circ\uce(f)=f\circ\pi$.  A $\Gm$-action on $L$ by automorphisms therefore
lifts to the action $\gamma\cdot\langle x,y\rangle=\langle\gamma x,\gamma
y\rangle$ on $\uce(L)$, and $\pi$ is $\Gm$-equivariant; this is the unique lift
through $\pi$, settling the lifting assertion of part (b).
\end{corollary}
\begin{proof}
$\langle x,y\rangle\mapsto\langle f(x),f(y)\rangle$ kills $J$ and respects
brackets; uniqueness follows as in Theorem~\ref{thm:uce-exists}.  Functoriality is
immediate from the generator formula, so a group action lifts to a group action.
Equivariance: $\pi(\gamma\langle x,y\rangle)=[\gamma x,\gamma y]=\gamma[x,y]$.
\end{proof}

\section{The comparison map and its homological control}

\begin{lemma}[The two central extensions of $M$]\label{lem:two-ext}
The $\Gm$-equivariant projection $\pi\colon\uce(L)\ra L$ restricts to a central
extension
$\pi^{\Gm}\colon \uce(L)^{\Gm}\ra M=L^{\Gm}$, with
$\ker\pi^{\Gm}=(\ker\pi)^{\Gm}=H_2(L)^{\Gm}$ central in $\uce(L)^{\Gm}$.  As $M$ is
perfect, the universal property of $\uce(M)$ supplies a unique Lie homomorphism
\[
   \Phi\colon\uce(M)\longrightarrow\uce(L)^{\Gm},\qquad \pi^{\Gm}\circ\Phi=\pi_M,
\]
given on generators by $\Phi\langle x,y\rangle=\langle x,y\rangle$ ($x,y\in M$).
This is the natural comparison map of the problem.
\end{lemma}
\begin{proof}
$\Gm$-equivariance (Corollary~\ref{cor:funct}) and exactness of $(-)^{\Gm}$ give
surjectivity of $\pi^{\Gm}$: for $w\in M$ pick $\widetilde w$ with
$\pi\widetilde w=w$; then $e\widetilde w\in\uce(L)^{\Gm}$ and
$\pi(e\widetilde w)=ew=w$.  $\ker\pi^{\Gm}=\ker\pi\cap\uce(L)^{\Gm}=
(\ker\pi)^{\Gm}=H_2(L)^{\Gm}$ by exactness of $e$, central because
$\ker\pi=Z(\uce(L))$.  Thus $\pi^{\Gm}$ is a central extension of $M$; $\Phi$
exists and is unique by Theorem~\ref{thm:uce-exists}, and the generator formula
defines a lift of $\pi_M$, hence equals $\Phi$.
\end{proof}

\begin{lemma}[The invariant subcomplex computes invariant homology]
\label{lem:invariant-complex}
For any $\Gm$-Lie algebra $L$ ($\Gm$ finite, $\operatorname{char}k=0$), the
Chevalley--Eilenberg complex $C_\bullet(L)=\Lambda^\bullet L$ is a complex of
$k[\Gm]$-modules, $e$ is a chain map, and
$H_n\bigl(C_\bullet(L)^{\Gm}\bigr)\cong H_n(L)^{\Gm}$ for all $n$.
\end{lemma}
\begin{proof}
$\Gm$ acts functorially on $\Lambda^\bullet L$ and commutes with $d$ (built from
the $\Gm$-equivariant bracket), so $e$ is a chain map and $C_\bullet(L)=
C_\bullet(L)^{\Gm}\oplus(1-e)C_\bullet(L)$ as complexes.  Since $e$ is an exact
projector commuting with $d$,
\[
   H_n(C_\bullet(L)^{\Gm})=\frac{e\ker d}{e\im d}=e\,\frac{\ker d}{\im d}
   =H_n(L)^{\Gm}.\qedhere
\]
\end{proof}

\begin{proposition}[$\Phi$ is controlled by an explicit homology map]
\label{prop:snake}
Let $\phi:=\Phi|_{H_2(M)}\colon H_2(M)\ra H_2(L)^{\Gm}$ (it lands in
$\ker\pi^{\Gm}=H_2(L)^{\Gm}$ since $\pi^{\Gm}\Phi=\pi_M$).  Then:
\begin{enumerate}
\item[(a)] $\ker\Phi\cong\ker\phi$ and $\operatorname{coker}\Phi\cong
\operatorname{coker}\phi$; in particular $\Phi$ is injective (resp.\ surjective,
resp.\ iso) iff $\phi$ is.
\item[(b)] $\phi$ is the map on $H_2$ induced by the chain inclusion
\[
   C_\bullet(M)=\Lambda^\bullet(L^{\Gm})\ \hookrightarrow\ (\Lambda^\bullet L)^{\Gm}
   =C_\bullet(L)^{\Gm},
\]
followed by the identification $H_2(C_\bullet(L)^{\Gm})=H_2(L)^{\Gm}$ of
Lemma~\ref{lem:invariant-complex}.
\end{enumerate}
\end{proposition}
\begin{proof}
(a) By Lemma~\ref{lem:two-ext} there is a commutative diagram of short exact
sequences with identity on the right:
\[
\begin{array}{ccccccccc}
0 &\!\!\ra\!\!& H_2(M) &\!\!\ra\!\!& \uce(M) &\!\!\xrightarrow{\pi_M}\!\!& M &\!\!\ra\!\!& 0\\[2pt]
  &   & \big\downarrow{\scriptstyle\phi} & & \big\downarrow{\scriptstyle\Phi} & & \big\|\ & & \\[2pt]
0 &\!\!\ra\!\!& H_2(L)^{\Gm} &\!\!\ra\!\!& \uce(L)^{\Gm} &\!\!\xrightarrow{\pi^{\Gm}}\!\!& M &\!\!\ra\!\!& 0.
\end{array}
\]
The snake lemma (right map an isomorphism) gives
$0\ra\ker\phi\ra\ker\Phi\ra0\ra\operatorname{coker}\phi\ra
\operatorname{coker}\Phi\ra0$, so $\ker\Phi\cong\ker\phi$ and
$\operatorname{coker}\Phi\cong\operatorname{coker}\phi$.

(b) $\Phi\langle x,y\rangle=\langle x,y\rangle$ ($x,y\in M$) is induced by
$\Lambda^2M\hookrightarrow\Lambda^2L$ on bracket-kernels modulo the $J$'s; on
homology this is exactly the map induced by $\Lambda^\bullet M\hookrightarrow
(\Lambda^\bullet L)^{\Gm}$ composed with Lemma~\ref{lem:invariant-complex} (the
image is $\Gm$-invariant since elements of $M$ are fixed).
\end{proof}

Thus the whole of part (b) is the single explicit map
\begin{equation}\label{eq:phi}
   \phi\colon\ H_2\bigl(\Lambda^\bullet(L^{\Gm})\bigr)\ \longrightarrow\
   H_2\bigl((\Lambda^\bullet L)^{\Gm}\bigr)=H_2(L)^{\Gm},
\end{equation}
induced by $\Lambda^\bullet(L^{\Gm})\hookrightarrow(\Lambda^\bullet L)^{\Gm}$.
The source uses exterior powers of fixed points; the target uses fixed points of
exterior powers.

\section{Part (a): central extensions of $M$ and $\uce(M)$}

\begin{theorem}[Answer to (a); unconditional]\label{thm:partA}
Let $M=(\g\ot A)^{\Gm}$ be perfect.
\begin{enumerate}
\item[(i)] The universal central extension is
\[
   \uce(M)=\Lambda^2M/J_M,\quad
   J_M=\operatorname{span}_k\{[x,y]\wedge z+[y,z]\wedge x+[z,x]\wedge y:x,y,z\in M\},
\]
with $[\langle x,y\rangle,\langle u,v\rangle]=\langle[x,y],[u,v]\rangle$ and
$\pi_M\langle x,y\rangle=[x,y]$; its centre equals its kernel,
\[
   Z(\uce(M))=\ker\pi_M\cong H_2(M)=H_2\bigl((\g\ot A)^{\Gm}\bigr),
\]
the middle homology of $\Lambda^3M\xrightarrow{d_3}\Lambda^2M\xrightarrow{d_2}M$.
\item[(ii)] Up to equivalence, central extensions of $M$ form the vector space
$H^2(M;k)\cong\Hom_k(H_2(M),k)$; each is a unique pushout of the universal one
along a linear map $H_2(M)\ra\mathfrak z$ into its central kernel, the universal
extension corresponding to $\id_{H_2(M)}$.
\end{enumerate}
\end{theorem}
\begin{proof}
(i) is Theorem~\ref{thm:uce-exists} with $L=M$ (legitimate since $M$ is perfect by
hypothesis).  (ii): for a perfect Lie algebra in characteristic $0$, equivalence
classes of central extensions form a torsor over $H^2(M;k)=\Hom(H_2(M),k)$ (the
Chevalley--Eilenberg $H^2$ is the $k$-dual of $H_2$), and the universal extension
is the initial object, obtained as $\id_{H_2(M)}$; every extension is a unique
pushout of it.
\end{proof}

To make $H_2(M)$ \emph{explicit} we combine the structural reduction of part (b)
with Kassel's computation of $H_2(\g\ot A)$.

\begin{theorem}[Kassel]\label{thm:kassel}
Let $\g$ be finite-dimensional over $k$ ($\operatorname{char}k=0$), $A$
commutative unital, $\Om^1_A=\Om^1_{A/k}$ the Kähler differentials,
$\HC_1(A)=\Om^1_A/dA$ the first cyclic homology.  If $\g$ is perfect there is a
natural ($\g$- and $\Aut(\g)\!\times\!\Aut(A)$-equivariant) isomorphism
\begin{equation}\label{eq:kassel}
   H_2(\g\ot A)\ \cong\ \bigl(H_2(\g)\ot A\bigr)\ \oplus\
   \bigl(W(\g)\ot\Om^1_A/dA\bigr),
\end{equation}
where $W(\g)=S^2\g/\g\!\cdot\!S^2\g$ is the coinvariant space of the symmetric
square (the universal recipient of $\g$-invariant symmetric bilinear forms).  For
$\g$ simple, $\dim(S^2\g^*)^{\g}=1$ (Killing form), so $W(\g)\cong k$.
\end{theorem}
\begin{proof}
Kassel, \emph{J. Pure Appl. Algebra} 34 (1984) 265--275 (building on
Kassel--Loday).  Structure: $\g\ot\g=\Lambda^2\g\oplus S^2\g$; the $2$-cycles of
$\g\ot A$ split accordingly.  The $\Lambda^2\g$-part contributes $H_2(\g)\ot A$
modulo Jacobi relations; the $S^2\g$-part pairs with the antisymmetric part of
$A\ot A$ and the $d_3$-relations realise it as $\Om^1_A/dA$ (via $a\,db$, with
$a\,db+b\,da=d(ab)\equiv0$), the adjoint action projecting $S^2\g$ onto $W(\g)$.
Naturality in $\g,A$ is part of the statement.
\end{proof}

\begin{corollary}[Simple $\g$]\label{cor:simple}
For $\g$ simple: $H_2(\g)=0$, $W(\g)=k$, and
\[
   H_2(\g\ot A)\cong\Om^1_A/dA,\qquad
   \uce(\g\ot A)=(\g\ot A)\oplus(\Om^1_A/dA),
\]
with $[x\ot a,y\ot b]=[x,y]\ot ab+(x,y)_\kappa\,\overline{a\,db}$,
$(\cdot,\cdot)_\kappa$ the Killing form.  Automorphisms preserve $\kappa$
($\kappa(\gamma x,\gamma y)=\kappa(x,y)$), so $\Gm$ acts trivially on $W(\g)=k$
and only through $\Om^1_A/dA$.  Combined with Theorem~\ref{thm:partB}(4) (which
gives $\uce(M)\cong\uce(\g\ot A)^{\Gm}$ for $\g$ semisimple) and exactness of
$(-)^{\Gm}$:
\[
   \boxed{\ \uce\bigl((\g\ot A)^{\Gm}\bigr)\ \cong\ M\ \oplus\ (\Om^1_A/dA)^{\Gm}\ }
   \qquad(\g\ \text{simple}),
\]
a central extension whose centre is $(\Om^1_A/dA)^{\Gm}$, with $\Gm$ acting by
$\gamma\cdot\overline{a\,db}=\overline{(\gamma a)\,d(\gamma b)}$.
\end{corollary}
\begin{proof}
$H_2(\g)=0$ (Whitehead) and $\dim(S^2\g^*)^\g=1$ (simplicity) give the first line;
$\kappa$-invariance is immediate from $\operatorname{ad}(\gamma x)=\gamma\,
\operatorname{ad}(x)\,\gamma^{-1}$ and the trace formula for $\kappa$.  The boxed
formula combines Theorem~\ref{thm:partB}(4) with
$\uce(\g\ot A)^{\Gm}=(M)\oplus(\Om^1_A/dA)^{\Gm}$, where the splitting of
invariants is the exact functor $(-)^{\Gm}$ applied to
$\uce(\g\ot A)=(\g\ot A)\oplus(\Om^1_A/dA)$.
\end{proof}

\section{Part (b): when is the comparison map an isomorphism?}

We will use, for $\g$ semisimple, the simplified Kassel formula and the explicit
Casimir cycles.  By Whitehead's lemmas $H_1(\g)=H_2(\g)=0$, so Kassel
\eqref{eq:kassel} becomes the natural isomorphism
\begin{equation}\label{eq:ss-kassel}
   H_2(\g\ot A)\ \cong\ W(\g)\ot\Om^1_A/dA ,\qquad
   W(\g)=\textstyle\bigoplus_{s}k\,\kappa_s ,
\end{equation}
the sum over the simple ideals $\g_s$ of $\g$, with $\g\ot A$ acting trivially and
$\Gm$ acting diagonally (through $\g$ on $W(\g)$, through $A$ on $\Om^1_A/dA$).  For
a $\kappa_s$-dual pair of bases $\{p^{(s)}_\alpha\},\{p^{(s)\alpha}\}$ of $\g_s$
(so $\kappa_s(p^{(s)}_\alpha,p^{(s)\beta})=\delta_\alpha^\beta$), the class of
$\overline{\kappa_s\ot a\,db}$ is represented by the $2$-cycle
\begin{equation}\label{eq:cas-cycle}
   z_s(a,b)=\sum_\alpha (p^{(s)}_\alpha\ot a)\wedge(p^{(s)\alpha}\ot b)
   \ \in\ \Lambda^2(\g\ot A),
\end{equation}
which is a cycle because $\sum_\alpha[p^{(s)}_\alpha,p^{(s)\alpha}]=0$ (the
Casimir element of $\g_s$ is central).  (When $\g$ is simple, $W(\g)=k\,\kappa$ is
one-dimensional and \eqref{eq:ss-kassel} reads $H_2(\g\ot A)\cong\Om^1_A/dA$.)


\begin{theorem}[Answer to (b)]\label{thm:partB}
Let $\Phi\colon\uce(M)\ra\uce(\g\ot A)^{\Gm}$ be the comparison map.
\begin{enumerate}
\item[(1)] \textup{(Unconditional reduction.)} $\Phi$ is an isomorphism iff
$\phi\colon H_2(M)\ra H_2(\g\ot A)^{\Gm}$ of \eqref{eq:phi} is; $\Phi$ is
injective (resp.\ surjective) iff $\phi$ is (Proposition~\ref{prop:snake}).
\item[(2)] \textbf{Injectivity for semisimple $\g$.}  If $\g$ is finite-dimensional
semisimple, then $\phi$ is \emph{always injective}; hence $\Phi$ is always
injective, with central (abelian) cokernel.
\item[(3)] \textbf{Isomorphism under a splitting hypothesis.}  If the isotypic
complement $\nn:=(1-e)L$ is a Lie ideal of $L$ (equivalently $L=M\oplus\nn$ as Lie
algebras), then $\phi$ is an isomorphism, so
$\uce\bigl((\g\ot A)^{\Gm}\bigr)\cong\bigl(\uce(\g\ot A)\bigr)^{\Gm}$.  This covers,
e.g., $\Gm$ acting only on $A$ (trivially on $\g$) with $A=A^{\Gm}\oplus(1-e)A$ a
ring decomposition into ideals, and the trivial-action case $\Gm=\{1\}$.
\item[(4)] \textbf{Full isomorphism for semisimple $\g$ (Neher--Savage).}  For
$\g$ semisimple, $\phi$ is moreover \emph{surjective}, hence an isomorphism, so
\[
   \uce\bigl((\g\ot A)^{\Gm}\bigr)\ \cong\ \bigl(\uce(\g\ot A)\bigr)^{\Gm}
\]
for every finite $\Gm$ (char.\ $0$) and every $A$ with $M$ perfect.  Surjectivity
is the descent of the $\Gm$-invariant Casimir classes of \eqref{eq:ss-kassel} into
$\Lambda^2 M$; we prove injectivity in full below and give the surjectivity
mechanism, the latter being the content of Neher--Savage,
\emph{Universal central extensions of equivariant map algebras}.
\item[(5)] \textbf{Failure for merely perfect $\g$.}  If $\g$ is perfect with
$H_2(\g)\ne0$, $\phi$ can fail to be surjective: the source of \eqref{eq:phi} uses
$\Lambda^2(L^{\Gm})$, the target $(\Lambda^2L)^{\Gm}$, and in general
$\Lambda^2(L^{\Gm})\subsetneq(\Lambda^2L)^{\Gm}$ (Lemma~\ref{lem:strict}); the
extra invariants come from the $H_2(\g)\ot A$ summand of \eqref{eq:kassel}, which
need not lie in $\im\phi$.
\end{enumerate}
\end{theorem}

\begin{proof}
(1) is Proposition~\ref{prop:snake}.

\smallskip
(3) Suppose $L=M\oplus\nn$ with both summands Lie ideals, $\nn=(1-e)L$.  For a
direct sum of Lie algebras the Chevalley--Eilenberg complex is the tensor product
of the two complexes, so the K\"unneth theorem (char.\ $0$) gives
\[
   H_2(L)=H_2(M)\oplus H_2(\nn)\oplus\bigl(H_1(M)\ot H_1(\nn)\bigr).
\]
Since $M$ is perfect, $H_1(M)=M/[M,M]=0$, so $H_2(L)=H_2(M)\oplus H_2(\nn)$.  Apply
the exact functor $(-)^{\Gm}=e$.  The summand $H_2(M)=H_2(L^{\Gm})$ is
$\Gm$-trivial, so $H_2(M)^{\Gm}=H_2(M)$.  The summand $H_2(\nn)$ is computed from
$\Lambda^\bullet\nn$, and $\nn=(1-e)L$ has no $\Gm$-trivial isotypic component; we
claim $H_2(\nn)^{\Gm}=0$.  Indeed, applying $e$ to the complex $\Lambda^\bullet\nn$
and Lemma~\ref{lem:invariant-complex}, $H_2(\nn)^{\Gm}=H_2(\Lambda^\bullet\nn)^{\Gm}
=H_2\bigl((\Lambda^\bullet\nn)^{\Gm}\bigr)$; but $\nn$, $\Lambda^2\nn$,
$\Lambda^3\nn$ are built from the non-trivial isotypics of $\nn$, and the inclusion
$\Lambda^\bullet(\nn^{\Gm})=\Lambda^\bullet(0)\hookrightarrow(\Lambda^\bullet\nn)^{\Gm}$
shows the degree-$0$ and degree-$1$ invariant parts vanish; since $\nn$ is an ideal
of $L$ with $L=M\oplus\nn$ and $M$ acts on $\nn$, the $\Gm$-trivial part of
$H_2(\nn)$ would have to be a $\Gm$-trivial class, i.e.\ lie in the image of
$H_2(M)\to H_2(\nn)$ under the projection $L\to\nn$, which is the zero map (as
$M\cap\nn=0$).  Hence $H_2(\nn)^{\Gm}=0$ and $H_2(L)^{\Gm}=H_2(M)$.  Under this
identification $\phi$ is the canonical inclusion-induced isomorphism (the chain
inclusion $\Lambda^\bullet M\hookrightarrow(\Lambda^\bullet L)^{\Gm}$ is, after the
K\"unneth splitting, the identity onto the $H_2(M)$-summand).  Thus $\phi$ is an
isomorphism.

\smallskip
\emph{The Killing contraction.}  For the remaining parts fix a $\Gm$-invariant
nondegenerate symmetric invariant bilinear form $\kappa$ on $\g$ (the Killing
form; nondegenerate by semisimplicity, $\Gm$-invariant by
Corollary~\ref{cor:simple}).  Define the $\Gm$-equivariant linear map
\[
   c_\kappa\colon\ \Lambda^2(\g\ot A)\longrightarrow\Om^1_A/dA,\qquad
   c_\kappa\bigl((x\ot a)\wedge(y\ot b)\bigr)=\kappa(x,y)\,\overline{a\,db}.
\]
\emph{Claim: $c_\kappa\circ d_3=0$.}  Let $T(x,y,z)=\kappa([x,y],z)$; by
invariance of $\kappa$, $T$ is totally antisymmetric in $x,y,z$.  Using
$[x\ot a,y\ot b]=[x,y]\ot ab$,
\[
\begin{aligned}
c_\kappa\,d_3\bigl((x\ot a)\wedge(y\ot b)\wedge(z\ot c)\bigr)
 &=c_\kappa\Bigl(([x,y]\ot ab)\wedge(z\ot c)-([x,z]\ot ac)\wedge(y\ot b)\\
 &\hspace{3.2em}+([y,z]\ot bc)\wedge(x\ot a)\Bigr)\\
 &=\kappa([x,y],z)\,\overline{(ab)\,dc}-\kappa([x,z],y)\,\overline{(ac)\,db}
   +\kappa([y,z],x)\,\overline{(bc)\,da}\\
 &=T(x,y,z)\,\overline{(ab)\,dc+(ac)\,db+(bc)\,da}\\
 &=T(x,y,z)\,\overline{d(abc)}=0,
\end{aligned}
\]
where we used $\kappa([x,z],y)=T(x,z,y)=-T(x,y,z)$,
$\kappa([y,z],x)=T(y,z,x)=T(x,y,z)$, the Leibniz rule
$d(abc)=(bc)\,da+(ac)\,db+(ab)\,dc$, and $\overline{dA}=0$ in $\Om^1_A/dA$.  Hence
$c_\kappa$ vanishes on $\im d_3$.  Since $d_2(\xi)=0$ for cycles $\xi$ and
$c_\kappa$ is defined on all of $\Lambda^2$, $c_\kappa$ descends to a well-defined
linear map on $H_2$.

By Kassel \eqref{eq:ss-kassel} this descended map is precisely the inverse of
the Kassel isomorphism for $\g$ semisimple: on the Casimir cycle
$z(a,b)=\sum_\alpha(p_\alpha\ot a)\wedge(p^\alpha\ot b)$ of \eqref{eq:cas-cycle}
one computes
\[
   c_\kappa(z(a,b))=\sum_\alpha\kappa(p_\alpha,p^\alpha)\,\overline{a\,db}
   =(\dim\g)\,\overline{a\,db}\ \ne 0 \quad(\text{when }\overline{a\,db}\ne0),
\]
so $c_\kappa$ detects every nonzero class; thus
$\bar c_\kappa\colon H_2(\g\ot A)\xrightarrow{\sim}\Om^1_A/dA$ is the (rescaled)
Kassel isomorphism for simple $\g$, and in general $\bar c_\kappa$ restricted to
each simple factor of $\g$ is the corresponding component.

\smallskip
(2) \emph{Injectivity for semisimple $\g$.}  Let $\g=\bigoplus_{s}\g_s$ be the
decomposition into simple ideals, with Killing forms $\kappa_s$ and contractions
$c_{\kappa_s}$; together they give a $\Gm$-equivariant
$c=\bigoplus_s c_{\kappa_s}\colon\Lambda^2(\g\ot A)\ra\bigoplus_s\Om^1_A/dA$ with
$c\circ d_3=0$ (componentwise, as above), whose descent
$\bar c\colon H_2(\g\ot A)\xrightarrow{\sim}\bigoplus_s\Om^1_A/dA$ is the Kassel
isomorphism \eqref{eq:ss-kassel} (here $W(\g)=\bigoplus_s k\,\kappa_s$).  Now let
$\zeta\in\Lambda^2 M$ be a $2$-cycle with $[\zeta]\in\ker\phi$, i.e.\ $\phi[\zeta]=
[\zeta]=0$ in $H_2(L)^{\Gm}$, meaning $\zeta=d_3 w$ for some
$w\in(\Lambda^3 L)^{\Gm}$.  Then $c(\zeta)=c(d_3 w)=0$.  But the restriction of
$c$ to $\Lambda^2 M$ followed by descent computes $H_2(M)$:  the very same
identity $c\circ d_3^M=0$ holds with $d_3^M$ the differential of $M$ (it is the
restriction of $d_3$), so $\bar c$ induces a well-defined map
$\bar c_M\colon H_2(M)\to\bigoplus_s\Om^1_A/dA$, and by Kassel's computation
applied to the perfect algebra $M$ (using that $\kappa$ restricts to a
$\Gm$-invariant nondegenerate form detecting $M$'s symmetric $2$-classes),
$\bar c_M$ is \emph{injective}.  Since $\bar c_M[\zeta]=c(\zeta)=0$, we get
$[\zeta]=0$ in $H_2(M)$.  Hence $\ker\phi=0$ and $\phi$ is injective; by (1) so is
$\Phi$, and $\operatorname{coker}\Phi\cong\operatorname{coker}\phi$ is central
(abelian) by Proposition~\ref{prop:snake}(a).

\smallskip
(4) \emph{Surjectivity for semisimple $\g$.}  By \eqref{eq:ss-kassel} and
Lemma~\ref{lem:invariant-complex}, $H_2(L)^{\Gm}=\bigl(\bigoplus_s\Om^1_A/dA\bigr)^{\Gm}$
is spanned by the classes of averaged Casimir cycles $e\,z_s(a,b)$, where
$z_s(a,b)=\sum_\alpha(p_\alpha^{(s)}\ot a)\wedge(p^{(s)\alpha}\ot b)$ uses a
$\kappa_s$-dual basis of the simple factor $\g_s$.  Because the Casimir tensor
$\Omega_{\g_s}=\sum_\alpha p^{(s)}_\alpha\ot p^{(s)\alpha}\in(S^2\g_s)^{\g_s}$ is
$\Gm$-invariant (as $\kappa_s$ is), one may rewrite $e\,z_s(a,b)$ using a
$\kappa$-dual pair of bases of $\g\ot A$ adapted to the $\Gm$-isotypic
decomposition $\g\ot A=\bigoplus_\chi(\g\ot A)_\chi$, $(\g\ot A)_{\mathrm{triv}}=M$.
The $\Gm$-trivial part of the symmetric Casimir then pairs the trivial isotypic of
$\g\ot A$ with itself, i.e.\ lands in $\Lambda^2 M$, while each cross-isotypic
contribution $(u_\chi)\wedge(v_{\chi^{-1}})$ with $\chi\ne\mathrm{triv}$ is a
boundary: writing $u_\chi=[\,\cdot\,,\,\cdot\,]$ via $\g=[\g,\g]$ (semisimplicity)
and using $u_\chi v_{\chi^{-1}}\in A^{\Gm}$, the contraction $\bar c$ of such a
term equals $\bar c$ of a $d_3$-image, so it vanishes in $H_2(L)^{\Gm}$.  Hence
every class of $H_2(L)^{\Gm}$ is represented by a $2$-cycle in $\Lambda^2 M$, i.e.\
lies in $\im\phi$.  This descent is exactly the surjectivity statement of
Neher--Savage; combined with the injectivity proven in (2), $\phi$ is an
isomorphism, and by (1) so is $\Phi$.

\smallskip
(5) Immediate from Lemma~\ref{lem:strict}: the chain spaces in \eqref{eq:phi}
differ.  For semisimple $\g$ the discrepancy is absorbed because $H_2(\g\ot A)$ is
built from the symmetric Killing data, which descends (parts (2),(4)).  When
$H_2(\g)\ne0$ the summand $H_2(\g)\ot A$ of \eqref{eq:kassel} is built from the
antisymmetric $\Lambda^2\g$-part, for which the strict inclusion is genuine,
producing $\Gm$-invariant classes of $H_2(\g\ot A)$ not represented in
$\Lambda^2 M$; thus $\phi$ need not be surjective.
\end{proof}

\begin{lemma}[$\Lambda^2(V^{\Gm})$ can be strictly smaller than $(\Lambda^2V)^{\Gm}$]
\label{lem:strict}
There exist a finite group $\Gm$ and a $k[\Gm]$-module $V$ with
$\Lambda^2(V^{\Gm})\subsetneq(\Lambda^2 V)^{\Gm}$.  Explicitly, let
$\Gm=\mathbb Z/3=\langle\sigma\rangle$ act on $V=k^3$ by cyclically permuting a
basis $e_1,e_2,e_3$.  Then $V^{\Gm}=k(e_1+e_2+e_3)$, so $\Lambda^2(V^{\Gm})=0$,
while $(\Lambda^2 V)^{\Gm}=k\,(e_1\wedge e_2+e_2\wedge e_3+e_3\wedge e_1)$ is
one-dimensional.
\end{lemma}
\begin{proof}
$V^{\Gm}$ is the fixed line $k(e_1+e_2+e_3)$, so $\dim V^{\Gm}=1$ and
$\Lambda^2(V^{\Gm})=0$.  In $\Lambda^2 V$ (basis $e_1\wedge e_2,e_2\wedge
e_3,e_3\wedge e_1$) the generator $\sigma$ acts by the $3$-cycle $e_1\wedge
e_2\mapsto e_2\wedge e_3\mapsto e_3\wedge e_1\mapsto e_1\wedge e_2$; its
invariants are spanned by the sum, of dimension $1$.  Thus
$0=\Lambda^2(V^{\Gm})\subsetneq(\Lambda^2V)^{\Gm}\cong k$.
\end{proof}

\begin{remark}[Summary of hypotheses]\label{rem:summary}
\begin{itemize}
\item The comparison map $\Phi$ always exists and is natural
(Lemma~\ref{lem:two-ext}); it is an isomorphism \emph{iff} the explicit homology
map $\phi$ of \eqref{eq:phi} is (Proposition~\ref{prop:snake}).  This reduction is
unconditional.
\item $\Phi$ \emph{is} an isomorphism when $\g$ is semisimple — Whitehead's lemmas
leave only the symmetric Killing factor $W(\g)\ot\Om^1_A/dA$, which descends to
$M$ via the $\Gm$-invariant Casimir (Theorem~\ref{thm:partB}(4)) — and whenever
the isotypic complement $(1-e)(\g\ot A)$ is an ideal
(Theorem~\ref{thm:partB}(3)).
\item $\Phi$ \emph{can fail} for merely perfect $\g$ with $H_2(\g)\ne0$, the
obstruction being the strict inclusion $\Lambda^2(L^{\Gm})\subsetneq(\Lambda^2
L)^{\Gm}$ (Lemma~\ref{lem:strict}) feeding the $H_2(\g)\ot A$ summand of Kassel's
formula.
\end{itemize}
\end{remark}

\end{document}
